% This is a new section to be added to hydrosheaf_technical_document.tex before the Conclusion section

\newpage
\section{Web-Based User Interface}\label{sec:webui}

To make Hydrosheaf accessible to a broader audience beyond command-line users, we have developed a modern full-stack web application (version 0.4.0). This interface provides interactive visualization, project management, and real-time analysis capabilities while maintaining the mathematical rigor of the underlying computational framework.

\subsection{Architecture Overview}

The web application follows a three-tier architecture:

\begin{itemize}
	\item \textbf{Frontend Layer}: React-based single-page application (SPA) with modern UI components, state management via React Context, and real-time updates via WebSockets.
	\item \textbf{Backend Layer}: FastAPI (Python 3.8+) REST API providing endpoints for project management, sample data validation, edge inference, and asynchronous analysis execution.
	\item \textbf{Computational Core}: The same mathematical framework presented in Sections \ref{sec:math_framework}--\ref{sec:extensions}, ensuring consistency between CLI and web interfaces. All transport optimization, LASSO solving, and thermodynamic constraints use identical implementations.
\end{itemize}

\subsection{Frontend Features}

The React-based frontend implements a "Dark Ocean" design theme with glassmorphism effects, providing an intuitive user experience for complex hydrogeochemical analysis:

\subsubsection{Interactive Dashboard}

The main dashboard provides:
\begin{itemize}
	\item \textbf{Project Management}: Create, load, and manage multiple groundwater analysis projects with persistent browser storage (IndexedDB) or server-side PostgreSQL backend.
	\item \textbf{Quick Statistics}: Real-time summaries of sample counts, inferred edges, and analysis completion status.
	\item \textbf{Recent Activity}: Timeline view of analysis runs, parameter changes, and results exports with version history.
\end{itemize}

\subsubsection{Sample Data Manager}

Provides spreadsheet-like interface for:
\begin{itemize}
	\item \textbf{Data Entry}: Input or paste chemical concentrations, coordinates, and metadata with automatic unit conversion (mg/L $\leftrightarrow$ mmol/L, meq/L $\leftrightarrow$ mmol/L).
	\item \textbf{CSV Import/Export}: Bulk upload of field data with schema validation (via Pydantic models) and error highlighting.
	\item \textbf{Quality Control}: Automated charge balance error calculation, outlier detection via robust Z-scores, and completeness validation (missing Ca, Mg, etc.).
\end{itemize}

\subsubsection{Network Visualization}

Interactive graph editor using D3.js force-directed layouts:
\begin{itemize}
	\item \textbf{Node Positioning}: Drag-and-drop arrangement with optional geographic coordinate alignment (lon/lat $\to$ Web Mercator projection).
	\item \textbf{Edge Inference Controls}: Visual configuration of probabilistic edge inference parameters (radius km, minimum head gradient, $p_{min}$ threshold).
	\item \textbf{3D Layer Support}: Toggle vertical stratification view and visualize cross-layer connectivity (requires \texttt{--3d} mode enabled).
	\item \textbf{Edge Attributes}: Color-coded edges by confidence probability, dominant transport model (evaporation/mixing), or total reaction extent $\|\mathbf{z}\|_1$.
\end{itemize}

\subsubsection{Analysis Configuration Panel}

Formalized interface for model parameters mapped to CLI flags (Table \ref{tab:impl_map}):
\begin{itemize}
	\item \textbf{Transport Models}: Select candidate endmembers (rainfall, seawater, river), enable/disable evaporation hypothesis, configure isotope penalty weights ($\eta_E$, $\eta_d$).
	\item \textbf{Reaction Dictionary}: Toggle minerals (calcite, gypsum, halite, silicates), set custom stoichiometry, select PHREEQC thermodynamic database (WATEQ4F, PHREEQC.dat, llnl.dat).
	\item \textbf{Regularization}: Adjust L1 penalty ($\lambda$), enable residence-time coupling (\texttt{--residence-coupling}), set coordinate descent convergence tolerance.
	\item \textbf{Advanced Options}: Enable temporal mode (\texttt{--temporal-enabled}), select uncertainty quantification method (Bootstrap/Bayesian MCMC), upload vadose priors CSV.
\end{itemize}

\subsubsection{Real-Time Analysis Monitor}

Asynchronous execution with live feedback via WebSockets (FastAPI's \texttt{WebSocketEndpoint}):
\begin{itemize}
	\item \textbf{Progress Tracking}: Per-edge status (queued, running, complete, failed) with estimated time remaining based on historical per-edge solve times.
	\item \textbf{Intermediate Results}: Streaming updates of transport model selection, LASSO iteration count, reaction convergence flags, and current objective value.
	\item \textbf{Log Viewer}: Real-time display of solver diagnostics, PHREEQC warnings (saturation index failures, database issues), and QC flags (charge balance errors, isotope inconsistencies).
\end{itemize}

\subsubsection{Results Dashboard}

Comprehensive visualization of analysis outputs:
\begin{itemize}
	\item \textbf{Summary Tables}: Sortable, filterable grids (AG Grid React) of edge results with dynamic columns: transport parameters ($\gamma$, $f$, endmember ID), reaction extents (one column per active mineral), fit metrics ($R^2$, residual norm, L1 norm).
	\item \textbf{Spatial Maps}: Leaflet-based choropleth visualization of transport model prevalence, TDS evolution gradients, nitrate source probabilities ($P_{\text{manure}}$).
	\item \textbf{Piper/Stiff Diagrams}: Automated hydrochemical facies classification using ternary projections and evolution trajectories along flow paths.
	\item \textbf{Time-Series Plots}: When temporal mode is enabled (\texttt{--temporal-enabled}), interactive Plotly charts of residence time estimates, seasonal decomposition (trend + periodic + residual), and TTD kernel visualizations (histogram of $g(\tau)$).
	\item \textbf{Uncertainty Bands}: Credible intervals (Bayesian MCMC) or bootstrap percentiles (BCa) overlaid on predicted concentrations when \texttt{--uncertainty} is enabled.
\end{itemize}

\subsection{Backend API}

The FastAPI backend exposes RESTful endpoints and WebSocket channels designed around resource-oriented patterns:

\subsubsection{Project Management Endpoints}

\begin{verbatim}
POST   /api/projects              Create new project
GET    /api/projects/{id}         Retrieve project metadata
PUT    /api/projects/{id}         Update configuration
DELETE /api/projects/{id}         Delete project
GET    /api/projects/{id}/export  Export as PDF report
\end{verbatim}

\subsubsection{Sample Data Endpoints}

\begin{verbatim}
POST /api/samples           Upload CSV or JSON sample data
GET  /api/samples/validate  Run schema and QC checks
PUT  /api/samples/{id}      Update individual sample
\end{verbatim}

\subsubsection{Security Architecture}

The backend implements several security measures to protect the integrity of scientific data and computational resources:

\begin{itemize}
	\item \textbf{Rate Limiting}: All API endpoints are protected by a token-bucket rate limiter (default 200 requests/minute) using \texttt{slowapi} to prevent abuse and denial-of-service attacks.
	\item \textbf{Input Sanitization}: Project identifiers and file paths are rigorously sanitized to prevent directory traversal attacks.
	\item \textbf{Validation Layer}: Pydantic schemas enforce strict type checking and range validation on all inputs before they reach the core computational engine.
	\item \textbf{Error Handling}: Production error handlers suppress stack traces to prevent information leakage while logging detailed diagnostics for administrators.
\end{itemize}

\subsubsection{Network Endpoints}


\begin{verbatim}
POST /api/edges/infer  Trigger probabilistic edge inference
GET  /api/edges         Retrieve inferred edges with p(i→j)
PUT  /api/edges/{id}    Manually override edge (enable/disable)
\end{verbatim}

\subsubsection{Analysis Endpoints}

\begin{verbatim}
POST /api/analysis/run            Start analysis (returns job_id)
WS   /api/analysis/ws/{job_id}   WebSocket for real-time progress
GET  /api/analysis/status/{job_id}  Poll job status
GET  /api/analysis/results/{job_id} Retrieve completed results
\end{verbatim}

\subsubsection{Data Transformation Layer}

The backend implements automatic conversions to maintain compatibility with the core framework:
\begin{itemize}
	\item \textbf{Unit Conversion}: Frontend accepts mg/L (common in water quality reports), backend converts to mmol/L using molar masses (e.g., $M_{\text{Ca}} = 40.08$ g/mol) before invoking \texttt{hydrosheaf.inference.network\_fit}.
	\item \textbf{Schema Validation}: Pydantic models enforce required fields (\texttt{site\_id}, \texttt{Ca}, \texttt{Mg}, \texttt{Na}, \texttt{HCO3}, \texttt{Cl}, \texttt{SO4}, \texttt{EC}, \texttt{TDS}, \texttt{pH}) and reject malformed inputs with HTTP 422 before reaching the solver.
	\item \textbf{Result Serialization}: Converts NumPy arrays (\texttt{ndarray}) and NetworkX graphs (\texttt{DiGraph}) to JSON-serializable formats (lists of dicts) for frontend consumption.
\end{itemize}

\subsection{Deployment and Scalability}

\subsubsection{Development Mode}

For local testing and demonstration:
\begin{verbatim}
# Terminal 1 (Backend)
cd web/backend
uvicorn app.main:app --reload --port 8000

# Terminal 2 (Frontend)
cd web/frontend
npm run dev  # Vite dev server on port 5173
\end{verbatim}

Open browser to \texttt{http://localhost:5173}. Frontend proxies API requests to \texttt{http://localhost:8000}.

\subsubsection{Production Deployment}

Recommended stack for institutional deployment (university research group, government agency):
\begin{itemize}
	\item \textbf{Reverse Proxy}: Nginx or Traefik for HTTPS termination (Let's Encrypt certificates), load balancing, and static asset serving.
	\item \textbf{Application Server}: Gunicorn with Uvicorn workers (4-8 workers for CPU-bound LASSO solving) for concurrent request handling.
	\item \textbf{Task Queue}: Celery + Redis for background LASSO/MCMC jobs exceeding HTTP timeout limits (default 30s for web APIs).
	\item \textbf{Database}: PostgreSQL 12+ for persistent project storage, user sessions, and analysis result caching. SQLite acceptable for single-user deployments.
	\item \textbf{File Storage}: S3-compatible object storage (MinIO, AWS S3) for large result exports (CSV/JSON $>$ 10 MB) and time-series data.
\end{itemize}

\subsubsection{Performance Considerations}

For networks with $> 100$ edges or Bayesian MCMC uncertainty quantification (4000+ samples per edge):
\begin{itemize}
	\item \textbf{Asynchronous Execution}: Long-running analyses (estimated time $> 30$s) execute in Celery background workers, preventing frontend timeout.
	\item \textbf{Incremental Results}: WebSocket streaming allows users to view partial results (e.g., first 10 edges completed) while remaining edges compute in parallel.
	\item \textbf{Caching}: Redis-backed result cache (TTL 1 hour) prevents redundant computation when users toggle visualization options (Piper vs. Stiff diagrams) without changing model parameters.
	\item \textbf{Database Indexing}: PostgreSQL B-tree indexes on \texttt{project\_id}, \texttt{job\_id}, and \texttt{edge\_id} enable sub-millisecond lookups for result retrieval.
\end{itemize}

\subsection{Integration with Existing Workflows}

The web interface is designed to complement, not replace, the CLI:

\begin{itemize}
	\item \textbf{Bidirectional Compatibility}: Projects created via web UI can be exported as CLI-compatible JSON configuration files (\texttt{hydrosheaf --config exported\_config.json}), and vice versa (CLI results can be imported into web UI for visualization).
	\item \textbf{Programmatic API}: Advanced users can bypass the web frontend and interact directly with the FastAPI backend via Python \texttt{requests} library or R \texttt{httr} package for automated workflows.
	\item \textbf{Batch Processing}: For large regional studies (hundreds of wells, dozens of time steps, MCMC uncertainty), the CLI remains preferred due to lower overhead and easier integration with HPC schedulers (SLURM, PBS).
\end{itemize}

\subsection{Accessibility and Usability}

The web interface follows modern accessibility best practices:
\begin{itemize}
	\item \textbf{WCAG 2.1 Level AA Compliance}: Sufficient color contrast ratios (4.5:1 for text), keyboard navigation support, ARIA labels for screen readers.
	\item \textbf{Responsive Design}: Mobile-friendly layouts (Tailwind CSS breakpoints) for tablet and phone access, though desktop recommended for complex network editing.
	\item \textbf{Internationalization (i18n)}: React-i18next framework supports multiple languages (English, Spanish, French planned for future releases).
\end{itemize}

\subsection{Future Development}

Planned enhancements to the web interface (see \texttt{web/INTEGRATION\_ROADMAP.md} for detailed timeline):

\begin{itemize}
	\item \textbf{Collaborative Features}: Multi-user access with role-based permissions (viewer, analyst, admin) using OAuth2/OIDC authentication (Keycloak, Auth0).
	\item \textbf{Interactive Tutorials}: Guided walkthroughs using synthetic datasets (similar to \texttt{hydrosheaf\_synthetic\_csv/}) to teach inverse modeling concepts interactively.
	\item \textbf{3D Visualization}: WebGL-based rendering (Three.js) of layered aquifer networks and vertical transport pathways, replacing 2D projections.
	\item \textbf{Parameter Sensitivity Analysis}: Interactive dashboards showing how $\lambda$, weights $\mathbf{W}$, and thermodynamic bounds affect sparsity $|\mathcal{A}|$ and fit quality $R^2$.
	\item \textbf{Integration with External Databases}: Connectors for USGS Water Quality Portal API, state monitoring databases (e.g., California GAMA), and remote PHREEQC thermodynamic data repositories.
	\item \textbf{Live Model Comparison}: Side-by-side comparison of multiple analysis runs (e.g., varying $\lambda$ from 0.01 to 1.0) with synchronized visualizations.
\end{itemize}

\subsection{Open-Source Licensing and Contributions}

The web application (both frontend and backend) is released under the same MIT License as the core Hydrosheaf package. Contributions welcome via GitHub pull requests. Frontend uses npm for dependency management (React 18, Vite 4, Tailwind CSS 3); backend uses Poetry for Python dependency management (FastAPI 0.104, Pydantic 2.5).

